\section{Glossary, notation, and reading guide}

This chapter collects the notation and terminology used in the paper
and in this manual, and gives reading routes for different audiences.

\subsection{Mathematical symbols}

\begin{center}
\small
\begin{tabular}{@{}p{0.20\textwidth}p{0.74\textwidth}@{}}
\toprule
Symbol & Meaning \\
\midrule
$T(x)$ & accelerated step $(5x+1)/2^{\vtwo{5x+1}}$ \\
$\Orbit_x(n)$ & $n$-th accelerated orbit value of $x$ \\
$t_j$ & exact step weight at depth $j$ \\
$\wordWeight_j$ & accumulated weight before step $j$ \\
$w$ & a word of step weights \\
$\wordA(w)$ & word molecule of $w$ \\
$m_j$ & margin $j\log_2 5-\wordWeight_j$ \\
$\alpha$ & $\log_2 5$ \\
$C_1,C_2,C_3$ & counts of weights $1$, $2$, at least $3$ \\
$g$ & orbit state $\Orbit_7(j-1)$ \\
$x$ & candidate predecessor \\
$r_j$ & block head \\
$r_s$ & terminal state \\
$B$ & block molecule \\
$u_i$ & segment step weights \\
$W_i$ & segment prefix weights \\
$H_s$ & unified-core threshold parameter \\
\bottomrule
\end{tabular}
\end{center}

\subsection{Reset-window symbols}

\begin{center}
\small
\begin{tabular}{@{}p{0.20\textwidth}p{0.74\textwidth}@{}}
\toprule
Symbol & Meaning \\
\midrule
$j$ & reset step depth \\
$k_0$ & exponent in $s5^{k_0}=r+1$ \\
$t$ & reset step weight, $t\in\{1,2\}$ \\
$\delta$ & reset parameter: $1$ for $t=1$, $\{1,3\}$ for $t=2$ \\
$s$ & odd terminal part with $5\nmid s$ \\
$r$ & previous terminal state \\
$n$ & tail length $s-j$ \\
$\Delta$ & weight difference $W_s-W_j$ \\
\bottomrule
\end{tabular}
\end{center}

\subsection{Status terms}

\begin{itemize}[leftmargin=*]
\item \textbf{Math: proven}: the mathematical argument is written out.
\item \textbf{Lean: compiled}: the Lean theorem compiles without
  \texttt{sorry}.
\item \textbf{Lean: pending}: the Lean theorem is not yet closed.
\item \textbf{compiled (conditional)}: the theorem compiles but
  consumes an open input.
\end{itemize}

\subsection{Reading routes}

\paragraph{For a mathematician.}
Read the main paper in order: definitions, framework, unified core,
bridge, Lean.  Use the supplementary chapters for the full derivations:
the block calculus chapter, the unified-core chapter, the CRT chapter,
and the divergence-bridge chapter.

\paragraph{For a Lean reader.}
Begin with the Lean theorem index and the reproducibility guide, then
read the statement shapes in the status matrix.  Follow the import
order: \texttt{WordWindow}, \texttt{Pmi}, \texttt{FinitePrefix},
\texttt{UnifiedCoreBridge}, \texttt{UnifiedCoreAudit},
\texttt{CycleBridge}, \texttt{DwdbDiv}.

\paragraph{For a referee.}
Check the pending ledger first.  Two statements are pending:
\texttt{unified\_\allowbreak core\_\allowbreak final\_\allowbreak
no\_\allowbreak hge} and
\texttt{five\_\allowbreak x\_\allowbreak plus\_\allowbreak one\_\allowbreak
diverges\_\allowbreak at\_\allowbreak 7}.  Then audit the dependency
graph to confirm that every compiled claim has a build.

\subsection{Cross-reference map}

\begin{center}
\small
\begin{tabular}{@{}p{0.34\textwidth}p{0.60\textwidth}@{}}
\toprule
Main-paper section & Supplementary chapters \\
\midrule
introduction and main theorem & glossary, status matrix \\
string-flow framework & vocabulary, PMI-B, block calculus \\
unified core & unified core, CRT, candidate exclusion \\
divergence bridge & cycle word algebra, window architecture,
  divergence bridge in full \\
Lean verification & Lean theorem index, reproducibility guide \\
appendix & all chapters \\
\bottomrule
\end{tabular}
\end{center}

\subsection{How to use the manual}

\begin{enumerate}[leftmargin=*]
\item Start with the glossary to fix notation.
\item Follow the reading route for your audience.
\item Use the status matrix for any Lean claim.
\item Use the dependency graph to find what a statement needs.
\item Never cite a pending statement as closed.
\end{enumerate}
